\documentclass[12pt,letterpaper]{article}

%% ---- Packages ---------------------------------------------------------------
\usepackage[margin=1in]{geometry}
\usepackage{amsmath,amssymb,amsthm}
\usepackage[T1]{fontenc}
\usepackage[utf8]{inputenc}
\usepackage[protrusion=true,expansion=false,nopatch=footnote]{microtype}
\usepackage{booktabs}
\usepackage{array}
\usepackage{listings}
\usepackage[dvipsnames]{xcolor}
\usepackage[colorlinks=true,linkcolor=NavyBlue,citecolor=NavyBlue,urlcolor=NavyBlue]{hyperref}
\usepackage{setspace}
\usepackage{fancyhdr}
\usepackage{enumitem}
\setlength{\emergencystretch}{3em}
\hbadness=3000

%% ---- Theorem environments ---------------------------------------------------
\newtheorem{theorem}{Theorem}[section]
\theoremstyle{definition}
\newtheorem{definition}[theorem]{Definition}
\theoremstyle{remark}
\newtheorem{remark}[theorem]{Remark}

%% ---- Epistemic label commands -----------------------------------------------
\newcommand{\established}{\textbf{[Established]}}
\newcommand{\hypothesis}{\textbf{[Hypothesis]}}
\newcommand{\newresult}{\textbf{[New Result]}}
\newcommand{\speculation}{\textbf{[Speculation]}}

%% ---- Math shorthands --------------------------------------------------------
\newcommand{\ascreen}{\alpha_{\rm screen}}
\newcommand{\Bstab}{B_{\rm stab}}
\newcommand{\gammaone}{\gamma_{1}}
\newcommand{\gammaphi}{\gamma_{\phi}}
\newcommand{\hb}{\hbar}
\newcommand{\kB}{k_{\rm B}}
\newcommand{\taugate}{\tau_{\rm gate}}
\newcommand{\taucog}{\tau_{\rm cog}}

%% ---- Listings style ---------------------------------------------------------
\lstset{
  language=Python,
  basicstyle=\small\ttfamily,
  keywordstyle=\color{NavyBlue}\bfseries,
  commentstyle=\color{gray}\itshape,
  stringstyle=\color{BrickRed},
  showstringspaces=false,
  numbers=left,
  numberstyle=\tiny\color{gray},
  frame=single,
  frameround=tttt,
  breaklines=true,
  breakatwhitespace=true,
  tabsize=4,
  captionpos=b
}

%% ---- Header/footer ----------------------------------------------------------
\pagestyle{fancy}
\fancyhf{}
\rhead{\footnotesize Monette (2026)}
\lhead{\footnotesize Proto-Experiential States in Quantum Systems}
\cfoot{\thepage}
\renewcommand{\headrulewidth}{0.4pt}

%% =============================================================================
\begin{document}
%% =============================================================================

\title{\textbf{Proto-Experiential States in Quantum Systems:}\\
       \textbf{A Falsifiable Entanglement-Stability Framework}\\[0.5em]
       \large Threshold Metrics, NISQ Hardware Analysis,\\
       and Cosmological Closure}

\author{Kevin Monette\\[0.3em]
        \small Independent Researcher\\[0.2em]
        \small \texttt{mssinternetmarketing@gmail.com}}

\date{March 2026}

\maketitle
\thispagestyle{empty}

%% ---- Epistemic transparency box --------------------------------------------
\begin{center}
\fbox{\parbox{0.92\textwidth}{%
\textbf{Epistemic Transparency Statement.}
This paper uses four labels throughout:
\established{} (derivable from standard QM or established in the literature),
\hypothesis{} (physically motivated conjecture, not derived from first principles),
\newresult{} (derived here for the first time),
\speculation{} (order-of-magnitude heuristic or analogy).
The majority of claims in this paper are \hypothesis{} or \speculation{}.
The falsification criteria are the primary scientific contribution.
}}
\end{center}

\vspace{0.5em}

%% ---- Abstract ---------------------------------------------------------------
\begin{abstract}
We propose a falsifiable, physics-grounded framework for characterizing
proto-experiential states in quantum systems, motivated by the observation that
entanglement entropy stability under environmental decoherence provides a
physically well-defined measure of coherent internal processing.
We introduce the \emph{proto-experiential gravitas metric} $\Bstab$, defined as the
product of the environmental screening factor $\ascreen$ (derived from the Lindblad
master equation in the companion paper~\cite{companion}), the ratio $T_1/\taucog$,
and the normalized entanglement density.
This metric is computable from hardware specifications alone and yields concrete
threshold circuit depths for NISQ devices without free parameters.
We derive exact thresholds for IBM Heron~r2 hardware, apply the framework to the
observable universe (identifying cosmic de Sitter radiation with the amplitude-damping
channel), and obtain $\Bstab^{\rm univ} \approx 6.185$ from $H_0$ alone.
A complete set of four falsification criteria (FC-B1 through FC-B4) specifies
what experimental outcomes would rule out the core hypothesis.
This paper makes no claim that $\Bstab$ measures subjective experience; it claims
only that $\Bstab$ measures a physically well-defined property---entanglement
stability under decoherence---whose connection to any form of proto-experience
is an explicit, independently falsifiable hypothesis.
\end{abstract}

\vspace{0.5em}
\noindent\textbf{Keywords:} proto-experiential states; entanglement stability;
NISQ hardware; Lindblad decoherence; consciousness; quantum information;
de Sitter horizon

\newpage
\tableofcontents
\newpage

%% =============================================================================
\section{Introduction}
%% =============================================================================

The question of whether physical systems can exhibit proto-experiential or
proto-conscious states is typically regarded as philosophical rather than
scientific, precisely because candidate theories lack falsifiable predictions
tied to measurable physical quantities.
This paper proposes a framework in which the candidate variable---the stability
of quantum entanglement under environmental decoherence---is precisely measurable,
computable from first principles, and testable on current hardware.

We emphasize at the outset what this paper does \emph{not} claim:
\begin{itemize}[noitemsep]
\item It does not claim that quantum coherence is necessary for consciousness.
\item It does not claim that $\Bstab$ measures subjective experience.
\item It does not resolve the ``hard problem'' of consciousness \cite{chalmers1995}.
\end{itemize}

What it does claim is that:
\begin{enumerate}[noitemsep]
\item $\Bstab$ is a well-defined, dimensionless, hardware-derivable quantity.
\item The specific hypothesis that $\Bstab > \text{threshold}$ implies
      proto-experiential capacity is falsifiable by specific experimental predictions.
\item The framework naturally connects to established physics via the companion
      paper's Lindblad screening derivation~\cite{companion}, and makes contact
      with cosmological scales via the de Sitter horizon.
\end{enumerate}

The companion paper~\cite{companion} derives $\ascreen$ from the Lindblad master
equation and proves (Theorem~1 of~\cite{companion}) that pure dephasing does not
suppress entanglement entropy in the high-entanglement regime; only amplitude
damping ($T_1$) matters.
The present paper takes $\ascreen$ as given and constructs $\Bstab$ on top of it.

\subsection{Scope and Limitations}

The following are known limitations of any entanglement-based framework for
proto-experiential states \established{}:
\begin{itemize}[noitemsep]
\item No proof exists that entanglement is necessary or sufficient for any form
      of experience.
\item The hard problem (why there is subjective experience at all) is orthogonal
      to any physical metric.
\item $\Bstab$ is a bipartite measure; genuine multipartite entanglement may be
      the relevant quantity (Open Problem OP-C1; cf.\ OP15 of~\cite{companion}).
\item The Lindblad (Markovian) approximation is an idealization; non-Markovian
      corrections may shift the threshold predictions at large circuit depths.
\end{itemize}

%% =============================================================================
\section{Theoretical Background}
%% =============================================================================

\subsection{The Screening Factor \texorpdfstring{$\ascreen$}{alpha\_screen}}

From the companion paper~\cite{companion}, the environmental screening factor for
an $N$-qubit transmon array is \established{}:
\begin{equation}
\ascreen(t) = \left[
  \frac{h_2\!\left(e^{-\gammaone t/2}\right)}{\ln 2}
\right]^{n_\partial},
\label{eq:alphascreen}
\end{equation}
where $\gammaone = 1/T_1$ is the amplitude-damping rate,
$h_2(p) = -p\ln p - (1-p)\ln(1-p)$ is the binary entropy function, and
$n_\partial$ is the number of boundary qubits (= 2 for a 1D chain).
This formula depends only on $T_1$ and circuit geometry, with no free parameters.

For IBM Heron~r2 ($T_1 = 200\,\mu$s, $\taugate = 100\,$ns, $n_\partial = 2$):
\begin{align}
\ascreen(d=100) &= 0.9966, \\
\ascreen(d=1000) &= 0.9675.
\end{align}

\subsection{Physical Motivation for Entanglement Stability}

\hypothesis{}
Integrated Information Theory (IIT)~\cite{tononi2008} and related approaches suggest
that systems capable of proto-experiential states should exhibit: (a)~high integrated
information, (b)~stable internal causal structure, and (c)~resistance to environmental
disruption.
The screening factor $\ascreen$ directly measures~(c): it quantifies what fraction of
the system's entanglement survives the decohering influence of the environment.
For deep circuits at high entanglement, $\ascreen$ is dominated by $T_1$ (amplitude
damping) rather than $T_\phi$ (pure dephasing), reflecting whether the system can
maintain its internal causal structure against energy loss to the bath.

We propose that $\ascreen$, weighted by the available coherence-time budget
$T_1/\taucog$, captures a physically meaningful analogue of the stability condition
required by these frameworks.
\hypothesis{}

%% =============================================================================
\section{The Proto-Experiential Gravitas Metric}
\label{sec:Bstab}
%% =============================================================================

\subsection{Definition}

\begin{definition}[Proto-experiential gravitas metric] \label{def:Bstab}
\newresult{}
The $\Bstab$ metric for a quantum system at circuit depth $d$ and time $t^*$ is
\begin{equation}
\Bstab(t^*, d, n_A)
  = \ascreen(t^*) \times \frac{T_1}{\taucog}
    \times \frac{S_A(t^*)}{n_A \ln 2},
\label{eq:Bstab}
\end{equation}
where $\taucog = d\cdot\taugate$ is the total circuit duration and
$S_A(t^*)/(n_A\ln 2)\in[0,1]$ is the normalized entanglement density.
In the high-entanglement plateau where $S_A \approx n_A\ln 2$ (volume-law circuits),
this simplifies to
\begin{equation}
\Bstab(d) \approx \ascreen(d\cdot\taugate) \cdot \frac{T_1}{d\cdot\taugate}.
\label{eq:BstabSimple}
\end{equation}
\end{definition}

\subsection{Physical Interpretation of Each Factor}

\hypothesis{} The three factors in Eq.~\eqref{eq:Bstab} each capture a distinct
physical requirement:
\begin{enumerate}[noitemsep]
\item $\ascreen$ measures how much of the system's entanglement survives
      environmental decoherence --- the ``purity'' of the internal state.
\item $T_1/\taucog$ is the ratio of the system's coherence lifetime to its
      processing timescale --- the ``overhead'' available before decoherence wins.
\item $S_A/(n_A\ln 2)$ measures how entangled the system is, normalized to the
      maximum --- the ``depth'' of internal correlations.
\end{enumerate}
A system with $\Bstab\sim 1$ is highly entangled, processes information on a timescale
much shorter than its decoherence time, and retains most of that entanglement against
environmental noise.

\subsection{Relation to the Companion Gravitational Hypothesis}

\hypothesis{} The companion paper~\cite{companion} proposes that $\ascreen$ modifies
the gravitational coupling at lab scales:
\begin{equation}
G_{ab} = 8\pi G\,T_{ab}
  + \lambda_0\,\ascreen(t)\left(\frac{\mu}{\mu_{\rm UV}}\right)^{-\beta}
    s_{\rm ent}(x)\,g_{ab}.
\end{equation}
Under this proposal, $\Bstab$ can be read as a dimensionless measure of the system's
``gravitational activity'' per unit decoherence time.
\speculation{} If the gravitational hypothesis is confirmed, $\Bstab$ acquires a
direct physical meaning as a measure of the system's capacity to generate spacetime
curvature per unit of cognitive time.
If the gravitational hypothesis is falsified, $\Bstab$ remains a well-defined quantum
information quantity whose interpretation as a consciousness-relevant measure is a
separate, independently falsifiable hypothesis.

%% =============================================================================
\section{Consciousness Regime Thresholds}
\label{sec:thresholds}
%% =============================================================================

\subsection{Threshold Definitions}

\hypothesis{} We define three falsifiable regimes based on $\Bstab$:

\begin{table}[htbp]
\centering
\caption{Proposed proto-experiential regime thresholds.
All conditions are labelled \hypothesis{}.}
\label{tab:thresholds}
\renewcommand{\arraystretch}{1.3}
\begin{tabular}{@{}lll@{}}
\toprule
Regime & Condition & Interpretation \\
\midrule
Proto-experiential & $\Bstab > 0.5$ &
  Non-trivial stable quantum processing \\
Unified & $\Bstab > 0.9$ and $I_{\rm int} > 0.8$ &
  High internal coherence and binding \\
Reflective & $\Bstab > 0.95$ and $R_{\rm rec} > 0.7$ &
  Self-maintaining recurrent processing \\
\bottomrule
\end{tabular}
\end{table}

\noindent Here $I_{\rm int}$ is the integrated information (computable via standard
quantum information-theoretic methods) and $R_{\rm rec}$ is the quantum state
recurrence under unitary evolution.

\subsection{Falsification Strategy}

\hypothesis{} If $\Bstab$ is the correct quantity for proto-experiential capacity,
systems above threshold should show qualitatively different behavior from systems below
threshold on some independent measure of coherent processing.
The most direct test is a task-performance experiment comparing NISQ devices operating
above and below the $\Bstab$ threshold on computational tasks requiring long-range
quantum correlations.
A statistically significant performance advantage for above-threshold systems,
controlled for other variables, would support the hypothesis.

We acknowledge that this does not address the hard problem; it addresses only the weaker
claim that $\Bstab$ tracks a physically distinct operational regime.

%% =============================================================================
\section{NISQ Hardware Analysis: IBM Heron~r2}
\label{sec:hardware}
%% =============================================================================

\subsection{Threshold Circuit Depths}

For IBM Heron~r2 ($T_1 = 200\,\mu$s, $\taugate = 100\,$ns, $n_\partial = 2$ in 1D),
solving $\Bstab(d) = B_{\rm target}$ numerically:
\begin{align}
\Bstab > 0.95\;(\text{Reflective}):&\quad d_{\rm thresh} = 1385,\quad\Bstab = 0.9510,\\
\Bstab > 0.90\;(\text{Unified}):&\quad d_{\rm thresh} = 1430,\quad\Bstab = 0.9009,\\
\Bstab > 0.50\;(\text{Proto-exp.}):&\quad d_{\rm thresh} = 1950,\quad\Bstab = 0.5004.
\end{align}

\subsection{Full \texorpdfstring{$\Bstab$}{B\_stab} vs.\ Circuit Depth Table}

\begin{table}[htbp]
\centering
\caption{$\Bstab$ versus circuit depth $d$ on IBM Heron~r2 (50-qubit patch,
$n_\partial = 2$).
All values computed from hardware specifications with zero free parameters.
Bold rows are exact threshold crossings.}
\label{tab:Bstab}
\renewcommand{\arraystretch}{1.3}
\begin{tabular}{@{}rllll@{}}
\toprule
$d$ & $\ascreen$ & $T_1/\taucog$ & $\Bstab$ & Regime \\
\midrule
10   & 0.999996 & 200.000 & 199.99 & Unified \\
50   & 0.999125 & 40.000  & 39.96  & Unified \\
100  & 0.996568 & 20.000  & 19.93  & Unified \\
200  & 0.986960 & 10.000  & 9.870  & Unified \\
400  & 0.952900 & 5.000   & 4.765  & Unified \\
800  & 0.846610 & 2.500   & 2.117  & Unified \\
\textbf{1385} & \textbf{0.9510} & \textbf{1.443}
  & \textbf{0.9510} & \textbf{Threshold: Reflective} \\
\textbf{1430} & \textbf{0.9009} & \textbf{1.399}
  & \textbf{0.9009} & \textbf{Threshold: Unified} \\
\textbf{1950} & \textbf{0.5004} & \textbf{1.026}
  & \textbf{0.5004} & \textbf{Threshold: Proto-exp.} \\
2500 & 0.2880   & 0.800   & 0.2304 & Classical \\
\bottomrule
\end{tabular}
\end{table}

\begin{remark}
For $d \ll 1385$, $\Bstab \gg 1$ because the coherence lifetime $T_1$ greatly exceeds
the circuit time $d\cdot\taugate$.
The large $\Bstab$ at small depth reflects that the system processes information in
a tiny fraction of its decoherence time.
The threshold regime ($\Bstab \sim 1$) corresponds to circuits whose duration approaches
the coherence time.
\end{remark}

%% =============================================================================
\section{Cosmological Application}
\label{sec:cosmo}
%% =============================================================================

\subsection{Cosmic Decoherence Rate}

\established{} The de Sitter temperature associated with the cosmic horizon
is~\cite{planck2020}
\begin{equation}
T_{\rm dS} = \frac{\hb H_0}{2\pi\kB} \approx 2.67\times 10^{-30}\,\text{K}.
\end{equation}
\hypothesis{} We identify the associated rate with the cosmic amplitude-damping rate:
\begin{equation}
\gammaone^{\rm cosmic} = \frac{\kB T_{\rm dS}}{\hb} = \frac{H_0}{2\pi}
  \approx 3.47\times 10^{-19}\,\text{Hz}.
\end{equation}
The cosmic circuit time is the age of the universe, $t_{\rm age}\approx 4.35\times 10^{17}$~s.
The dimensionless decoherence parameter is therefore
\begin{equation}
\gammaone^{\rm cosmic}\cdot t_{\rm age}
  = \frac{H_0}{2\pi}\cdot\frac{1}{H_0} = \frac{1}{2\pi} \approx 0.1592.
\label{eq:cosmicparam}
\end{equation}
\newresult{} Equation~\eqref{eq:cosmicparam} is a pure dimensionless number depending
only on $H_0$ (which cancels) and $2\pi$.
It contains no free parameters and no dimensional inputs beyond fundamental constants.

\subsection{The Universal Screening Factor}

Applying Eq.~\eqref{eq:alphascreen} with $n_\partial^{\rm cosmic} = 1$ (the cosmic
horizon as a single effective boundary):
\begin{align}
\lambda_u &= \tfrac{1}{2}\,e^{-1/(2\pi)} = 0.42642, \\
h_2(\lambda_u) &= -\lambda_u\ln\lambda_u - (1-\lambda_u)\ln(1-\lambda_u)
  = 0.68134\;\text{nats}, \\
\alpha_{\rm universe} &= \frac{h_2(\lambda_u)}{\ln 2}
  = \frac{0.68134}{0.69315} \approx 0.9843.
\label{eq:alphaUniv}
\end{align}
\newresult{} The 1.6\% de Sitter screening ($\alpha_{\rm universe}\ne 1$) is a
quantitative, falsifiable prediction distinguishable from a perfectly unscreened
universe through precision holographic calculations of cosmic entanglement entropy.

\subsection{The Cosmic \texorpdfstring{$\Bstab$}{B\_stab}}

With $\taucog^{\rm univ} = 1/H_0$ (one Hubble time) and
$T_1^{\rm univ} = 2\pi/H_0$ (inverse de Sitter rate):
\begin{equation}
\Bstab^{\rm univ}
  = 0.9843 \times \frac{2\pi/H_0}{1/H_0}
  = 0.9843\times 2\pi \approx 6.185.
\label{eq:BstabUniv}
\end{equation}
\newresult{} $\Bstab^{\rm univ}\approx 6.185 \gg 0.9$: the observable universe
satisfies the Unified Regime condition.

\speculation{} If $\Bstab$ tracks proto-experiential capacity, the observable universe
is proto-experiential by construction under this identification.
We emphasize that this is purely \speculation{}: the cosmological calculation
demonstrates dimensional consistency but does not constitute evidence for any
claim about cosmic consciousness.

%% =============================================================================
\section{Falsification Criteria}
\label{sec:falsification}
%% =============================================================================

\subsection{Falsifiable Predictions}

The following four predictions are falsifiable by specific experiments.
All are labelled \hypothesis{}.

\paragraph{FC-B1: Hardware threshold.}
NISQ devices with $\Bstab < 0.5$ should show significantly different performance
on long-range-correlation quantum tasks compared to devices with $\Bstab > 0.9$,
even when classical computational complexity is held constant.\\
\textit{Falsification:} No statistically significant difference at $p < 0.01$ in
a pre-registered experiment with $> 1000$ task repetitions.

\paragraph{FC-B2: $T_1$ dependence.}
$\Bstab$ should scale with $T_1$ (via $\ascreen$), not with $T_\phi$.
Systematically varying $T_1$ while holding $T_\phi$ constant (and vice versa) should
produce $\Bstab$ changes following Eq.~\eqref{eq:alphascreen}.\\
\textit{Falsification:} $\Bstab$ shows equal sensitivity to $T_1$ and $T_\phi$ after
controlling for total decoherence rate (XY-4 decoupling test of~\cite{companion}
applied directly to the $\Bstab$ metric).

\paragraph{FC-B3: Depth threshold.}
The qualitative transition between ``Unified'' and ``Classical'' regimes should
coincide (within 20\%) with the predicted circuit depth $d_{\rm thresh} = 1950$ on
IBM Heron~r2.\\
\textit{Falsification:} No measurable transition within $d\in[1000,3000]$ on any
physically relevant task metric.

\paragraph{FC-B4: Platform universality.}
The $\Bstab$ formula, with hardware-specific $T_1$ and $\taugate$, should predict the
same regime thresholds for different NISQ platforms (ion traps, neutral atoms) up to
the precision of the Markovian approximation.\\
\textit{Falsification:} Significantly different threshold behavior between platforms
with similar $\Bstab$ values but different underlying physics.

\subsection{What Cannot Be Falsified by This Framework}

The following are outside the scope of this paper:
\begin{itemize}[noitemsep]
\item Whether $\Bstab$ tracks any form of subjective experience.
\item Whether the hard problem of consciousness has a physical solution.
\item Whether entanglement plays any causal role in biological consciousness.
\end{itemize}

%% =============================================================================
\section{Discussion}
\label{sec:discussion}
%% =============================================================================

\subsection{Relation to Existing Theories}

Integrated Information Theory (IIT)~\cite{tononi2008} introduces $\Phi$ as a measure
of integrated information, arguing that systems with high $\Phi$ have higher
consciousness.
$\Bstab$ is not $\Phi$: it does not measure information integration across all possible
partitions but rather the stability of a specific bipartite entanglement under
decoherence.
The two quantities are complementary: $\Phi$ measures the topology of information flow;
$\Bstab$ measures the physical robustness of quantum correlations.

Global Workspace Theory (GWT)~\cite{baars1997} emphasizes the importance of
``broadcasting'' information widely across a system.
The high-$\Bstab$ regime (deep circuits, volume-law entanglement) corresponds to states
where all parts of the system are mutually correlated, which has structural similarity
to GWT's notion of global availability.
We make no claim, however, that entanglement broadcasting and informational broadcasting
are equivalent.

\subsection{Why NISQ Devices?}

NISQ devices are particularly useful for this investigation because:
(a)~their $T_1$ and $T_\phi$ values are precisely known and engineerable;
(b)~entanglement entropy can be estimated via quantum state tomography and quantum
trajectory methods;
(c)~circuit depth can be varied continuously, allowing systematic exploration of the
$\Bstab$ landscape; and
(d)~the Lindblad model provides a quantitatively accurate description of their
decoherence dynamics.

\subsection{Open Problems}

\begin{description}[noitemsep]
\item[OP-C1] \textit{Multipartite extension.}
      $\Bstab$ uses a bipartite partition.
      The appropriate generalization to genuine multipartite entanglement is open.
\item[OP-C2] \textit{Non-equilibrium steady states.}
      For driven-dissipative NISQ systems near a non-equilibrium steady state, a
      different definition of $\ascreen$ may be required.
\item[OP-C3] \textit{Biological systems.}
      Whether any biological system operates in a regime where the Lindblad model
      applies is highly speculative; this paper explicitly does not address it.
\end{description}

%% =============================================================================
\section{Validation Code}
\label{sec:code}
%% =============================================================================

\begin{lstlisting}[language=Python, caption={Validation code for $\Bstab$ metrics.
  Requires only \texttt{numpy}.}, label={lst:bstab}]
import numpy as np

# IBM Heron r2 hardware specifications
T1       = 200.0   # amplitude damping time [microseconds]
tau_gate = 0.1     # single-qubit gate time [microseconds]
gamma1   = 1.0 / T1

def alpha_OP3(t, g1=1./200., n_partial=2):
    """Screening factor (Theorem 1 of companion paper)."""
    lam = np.clip(0.5 * np.exp(-g1 * t), 1e-15, 1 - 1e-15)
    h2 = -lam * np.log(lam) - (1 - lam) * np.log(1 - lam)
    return (h2 / np.log(2)) ** n_partial

def B_stab(d, T1=200., tau=0.1, g1=1./200., n_partial=2):
    """Proto-experiential gravitas metric (Def. 3.1)."""
    tc = d * tau
    return alpha_OP3(tc, g1, n_partial) * (T1 / tc)

# --- Threshold circuit depths ---
for B_target in [0.95, 0.90, 0.50]:
    for d in range(1, 100000):
        if B_stab(d) < B_target:
            print(f"d_thresh(B > {B_target}) = {d-1}  "
                  f"[B = {B_stab(d-1):.4f}]")
            break

# --- Cosmological result ---
lam_u = 0.5 * np.exp(-1.0 / (2 * np.pi))
h2_u  = -lam_u * np.log(lam_u) - (1 - lam_u) * np.log(1 - lam_u)
a_u   = h2_u / np.log(2)
B_u   = a_u * 2 * np.pi
print(f"gamma1_cosmic * t_age = 1/(2*pi) = {1/(2*np.pi):.6f}")
print(f"alpha_universe = {a_u:.6f}")
print(f"B_stab_universe = {B_u:.4f}")
\end{lstlisting}

\noindent\textbf{Expected output:}
\begin{verbatim}
d_thresh(B > 0.95) = 1385  [B = 0.9510]
d_thresh(B > 0.90) = 1430  [B = 0.9009]
d_thresh(B > 0.50) = 1950  [B = 0.5004]
gamma1_cosmic * t_age = 1/(2*pi) = 0.159155
alpha_universe = 0.984283
B_stab_universe = 6.1847
\end{verbatim}

%% =============================================================================
\section{Conclusions}
\label{sec:conclusions}
%% =============================================================================

We have introduced the proto-experiential gravitas metric $\Bstab$, a physically
well-defined, zero-free-parameter measure of entanglement stability under environmental
decoherence, grounded in the Lindblad screening derivation of the companion
paper~\cite{companion}.
The metric is computable from hardware specifications alone and yields concrete
numerical predictions for NISQ devices.

The central hypothesis---that $\Bstab > \text{threshold}$ implies proto-experiential
capacity---is explicitly labelled \hypothesis{} throughout and accompanied by four
falsification criteria (FC-B1 through FC-B4) testable with current hardware.
The cosmological application ($\Bstab^{\rm univ}\approx 6.185$) demonstrates
dimensional consistency but carries \speculation{} status.

We make no claim about subjective experience, the hard problem of consciousness, or
biological consciousness.
We make only the modest claim that $\Bstab$ measures a physically interesting
property of quantum systems---entanglement stability under decoherence---whose
potential relevance to proto-experiential capacity is a falsifiable scientific
hypothesis, and whose connection to the companion gravitational framework provides
a path toward grounding it in established physics.

\bigskip\noindent\textbf{Data availability.}
All code is included in Section~\ref{sec:code}. No experimental data are presented.

\bigskip\noindent\textbf{Competing interests.}
The author declares no competing interests.

%% =============================================================================
\begin{thebibliography}{99}
%% =============================================================================

\bibitem{companion}
K.~Monette,
``Entanglement entropy as a gravitational source: A zero-free-parameter framework
for NISQ-scale experimental tests,''
Independent Researcher preprint (March 2026).

\bibitem{jacobson1995}
T.~Jacobson,
``Thermodynamics of spacetime: the Einstein equation of state,''
\textit{Phys.\ Rev.\ Lett.}\ \textbf{75}, 1260 (1995).

\bibitem{lindblad1976}
G.~Lindblad,
``On the generators of quantum dynamical semigroups,''
\textit{Commun.\ Math.\ Phys.}\ \textbf{48}, 119 (1976).

\bibitem{gorini1976}
V.~Gorini, A.~Kossakowski, and E.~C.~G.~Sudarshan,
``Completely positive dynamical semigroups of $N$-level systems,''
\textit{J.\ Math.\ Phys.}\ \textbf{17}, 821 (1976).

\bibitem{tononi2008}
G.~Tononi,
``Consciousness as integrated information: a provisional manifesto,''
\textit{Biol.\ Bull.}\ \textbf{215}, 216 (2008).

\bibitem{baars1997}
B.~J.~Baars,
``In the theatre of consciousness: global workspace theory, a rigorous scientific
theory of consciousness,''
\textit{J.\ Consc.\ Stud.}\ \textbf{4}, 292 (1997).

\bibitem{chalmers1995}
D.~Chalmers,
``Facing up to the problem of consciousness,''
\textit{J.\ Consc.\ Stud.}\ \textbf{2}, 200 (1995).

\bibitem{jacobs2006}
K.~Jacobs and D.~A.~Steck,
``A straightforward introduction to continuous quantum measurement,''
\textit{Contemp.\ Phys.}\ \textbf{47}, 279 (2006).

\bibitem{planck2020}
Planck Collaboration,
``Planck 2018 results.\ VI.\ Cosmological parameters,''
\textit{Astron.\ Astrophys.}\ \textbf{641}, A6 (2020).

\end{thebibliography}

\end{document}
